%Diagrama de casos de uso.
%Se incluye vía %Diagrama de casos de uso.
%Se incluye vía %Diagrama de casos de uso.
%Se incluye vía %Diagrama de casos de uso.
%Se incluye vía \input{casos_uso} desde secciones/06_desarrollo.tex.
\begin{tikzpicture}[
    every node/.style={font=\small},
    actor/.style={
        draw, thick, minimum height=1cm, minimum width=1cm,
        align=center, font=\scriptsize
    },
    caso/.style={
        ellipse, draw, thick, align=center, font=\scriptsize,
        minimum height=0.75cm, minimum width=2.6cm, fill=black!3
    },
    conn/.style={-, thick},
    sistema/.style={
        draw, thick, dashed, rounded corners=6pt, inner sep=12pt
    }
]

%Actores
\node[actor] (docente) at (-6.5, 0)
    {\begin{tabular}{c}$\dagger$\\ Docente\end{tabular}};
\node[actor] (admin) at (-6.5, -5)
    {\begin{tabular}{c}$\dagger$\\ Administrador\end{tabular}};
\node[actor] (publico) at (-6.5, -9.5)
    {\begin{tabular}{c}$\dagger$\\ Público\end{tabular}};

%Casos de uso del docente
\node[caso] (login1) at (0, 2)      {Iniciar sesión};
\node[caso] (uc1)    at (0, 1)      {Gestionar cartas\\temáticas};
\node[caso] (uc2)    at (0, 0)      {Gestionar requisitos\\de recuperación};
\node[caso] (uc3)    at (0, -1)     {Responder\\autoevaluación};
\node[caso] (uc4)    at (0, -2)     {Consultar propio\\perfil e historial};

%Casos de uso del administrador
\node[caso] (uc5)   at (0, -3.5)  {Gestionar usuarios\\y profesores};
\node[caso] (uc6)   at (0, -4.5)  {Administrar\\catálogos};
\node[caso] (uc7)   at (0, -5.5)  {Diseñar formularios\\de autoevaluación};
\node[caso] (uc8)   at (0, -6.5)  {Consultar cartas y\\requisitos de todos};
\node[caso] (uc9)   at (0, -7.5)  {Generar reportes\\de cumplimiento};

%Casos de uso del público
\node[caso] (uc10)  at (0, -9)    {Explorar UEA por\\periodo y programa};
\node[caso] (uc11)  at (0, -10)   {Ver carta o requisito\\publicado};

%Sistema
\begin{scope}[on background layer]
\node[sistema, fit=(login1)(uc1)(uc2)(uc3)(uc4)(uc5)(uc6)(uc7)(uc8)(uc9)(uc10)(uc11)] (sys) {};
\end{scope}
\node at (0, 2.9) {\textbf{Sistema CyAD}};

%Conexiones
\draw[conn] (docente) -- (login1);
\draw[conn] (docente) -- (uc1);
\draw[conn] (docente) -- (uc2);
\draw[conn] (docente) -- (uc3);
\draw[conn] (docente) -- (uc4);

\draw[conn] (admin) -- (login1);
\draw[conn] (admin) -- (uc5);
\draw[conn] (admin) -- (uc6);
\draw[conn] (admin) -- (uc7);
\draw[conn] (admin) -- (uc8);
\draw[conn] (admin) -- (uc9);

\draw[conn] (publico) -- (uc10);
\draw[conn] (publico) -- (uc11);

\end{tikzpicture}
 desde secciones/06_desarrollo.tex.
\begin{tikzpicture}[
    every node/.style={font=\small},
    actor/.style={
        draw, thick, minimum height=1cm, minimum width=1cm,
        align=center, font=\scriptsize
    },
    caso/.style={
        ellipse, draw, thick, align=center, font=\scriptsize,
        minimum height=0.75cm, minimum width=2.6cm, fill=black!3
    },
    conn/.style={-, thick},
    sistema/.style={
        draw, thick, dashed, rounded corners=6pt, inner sep=12pt
    }
]

%Actores
\node[actor] (docente) at (-6.5, 0)
    {\begin{tabular}{c}$\dagger$\\ Docente\end{tabular}};
\node[actor] (admin) at (-6.5, -5)
    {\begin{tabular}{c}$\dagger$\\ Administrador\end{tabular}};
\node[actor] (publico) at (-6.5, -9.5)
    {\begin{tabular}{c}$\dagger$\\ Público\end{tabular}};

%Casos de uso del docente
\node[caso] (login1) at (0, 2)      {Iniciar sesión};
\node[caso] (uc1)    at (0, 1)      {Gestionar cartas\\temáticas};
\node[caso] (uc2)    at (0, 0)      {Gestionar requisitos\\de recuperación};
\node[caso] (uc3)    at (0, -1)     {Responder\\autoevaluación};
\node[caso] (uc4)    at (0, -2)     {Consultar propio\\perfil e historial};

%Casos de uso del administrador
\node[caso] (uc5)   at (0, -3.5)  {Gestionar usuarios\\y profesores};
\node[caso] (uc6)   at (0, -4.5)  {Administrar\\catálogos};
\node[caso] (uc7)   at (0, -5.5)  {Diseñar formularios\\de autoevaluación};
\node[caso] (uc8)   at (0, -6.5)  {Consultar cartas y\\requisitos de todos};
\node[caso] (uc9)   at (0, -7.5)  {Generar reportes\\de cumplimiento};

%Casos de uso del público
\node[caso] (uc10)  at (0, -9)    {Explorar UEA por\\periodo y programa};
\node[caso] (uc11)  at (0, -10)   {Ver carta o requisito\\publicado};

%Sistema
\begin{scope}[on background layer]
\node[sistema, fit=(login1)(uc1)(uc2)(uc3)(uc4)(uc5)(uc6)(uc7)(uc8)(uc9)(uc10)(uc11)] (sys) {};
\end{scope}
\node at (0, 2.9) {\textbf{Sistema CyAD}};

%Conexiones
\draw[conn] (docente) -- (login1);
\draw[conn] (docente) -- (uc1);
\draw[conn] (docente) -- (uc2);
\draw[conn] (docente) -- (uc3);
\draw[conn] (docente) -- (uc4);

\draw[conn] (admin) -- (login1);
\draw[conn] (admin) -- (uc5);
\draw[conn] (admin) -- (uc6);
\draw[conn] (admin) -- (uc7);
\draw[conn] (admin) -- (uc8);
\draw[conn] (admin) -- (uc9);

\draw[conn] (publico) -- (uc10);
\draw[conn] (publico) -- (uc11);

\end{tikzpicture}
 desde secciones/06_desarrollo.tex.
\begin{tikzpicture}[
    every node/.style={font=\small},
    actor/.style={
        draw, thick, minimum height=1cm, minimum width=1cm,
        align=center, font=\scriptsize
    },
    caso/.style={
        ellipse, draw, thick, align=center, font=\scriptsize,
        minimum height=0.75cm, minimum width=2.6cm, fill=black!3
    },
    conn/.style={-, thick},
    sistema/.style={
        draw, thick, dashed, rounded corners=6pt, inner sep=12pt
    }
]

%Actores
\node[actor] (docente) at (-6.5, 0)
    {\begin{tabular}{c}$\dagger$\\ Docente\end{tabular}};
\node[actor] (admin) at (-6.5, -5)
    {\begin{tabular}{c}$\dagger$\\ Administrador\end{tabular}};
\node[actor] (publico) at (-6.5, -9.5)
    {\begin{tabular}{c}$\dagger$\\ Público\end{tabular}};

%Casos de uso del docente
\node[caso] (login1) at (0, 2)      {Iniciar sesión};
\node[caso] (uc1)    at (0, 1)      {Gestionar cartas\\temáticas};
\node[caso] (uc2)    at (0, 0)      {Gestionar requisitos\\de recuperación};
\node[caso] (uc3)    at (0, -1)     {Responder\\autoevaluación};
\node[caso] (uc4)    at (0, -2)     {Consultar propio\\perfil e historial};

%Casos de uso del administrador
\node[caso] (uc5)   at (0, -3.5)  {Gestionar usuarios\\y profesores};
\node[caso] (uc6)   at (0, -4.5)  {Administrar\\catálogos};
\node[caso] (uc7)   at (0, -5.5)  {Diseñar formularios\\de autoevaluación};
\node[caso] (uc8)   at (0, -6.5)  {Consultar cartas y\\requisitos de todos};
\node[caso] (uc9)   at (0, -7.5)  {Generar reportes\\de cumplimiento};

%Casos de uso del público
\node[caso] (uc10)  at (0, -9)    {Explorar UEA por\\periodo y programa};
\node[caso] (uc11)  at (0, -10)   {Ver carta o requisito\\publicado};

%Sistema
\begin{scope}[on background layer]
\node[sistema, fit=(login1)(uc1)(uc2)(uc3)(uc4)(uc5)(uc6)(uc7)(uc8)(uc9)(uc10)(uc11)] (sys) {};
\end{scope}
\node at (0, 2.9) {\textbf{Sistema CyAD}};

%Conexiones
\draw[conn] (docente) -- (login1);
\draw[conn] (docente) -- (uc1);
\draw[conn] (docente) -- (uc2);
\draw[conn] (docente) -- (uc3);
\draw[conn] (docente) -- (uc4);

\draw[conn] (admin) -- (login1);
\draw[conn] (admin) -- (uc5);
\draw[conn] (admin) -- (uc6);
\draw[conn] (admin) -- (uc7);
\draw[conn] (admin) -- (uc8);
\draw[conn] (admin) -- (uc9);

\draw[conn] (publico) -- (uc10);
\draw[conn] (publico) -- (uc11);

\end{tikzpicture}
 desde secciones/06_desarrollo.tex.
\begin{tikzpicture}[
    every node/.style={font=\small},
    actor/.style={
        draw, thick, minimum height=1cm, minimum width=1cm,
        align=center, font=\scriptsize
    },
    caso/.style={
        ellipse, draw, thick, align=center, font=\scriptsize,
        minimum height=0.75cm, minimum width=2.6cm, fill=black!3
    },
    conn/.style={-, thick},
    sistema/.style={
        draw, thick, dashed, rounded corners=6pt, inner sep=12pt
    }
]

%Actores
\node[actor] (docente) at (-6.5, 0)
    {\begin{tabular}{c}$\dagger$\\ Docente\end{tabular}};
\node[actor] (admin) at (-6.5, -5)
    {\begin{tabular}{c}$\dagger$\\ Administrador\end{tabular}};
\node[actor] (publico) at (-6.5, -9.5)
    {\begin{tabular}{c}$\dagger$\\ Público\end{tabular}};

%Casos de uso del docente
\node[caso] (login1) at (0, 2)      {Iniciar sesión};
\node[caso] (uc1)    at (0, 1)      {Gestionar cartas\\temáticas};
\node[caso] (uc2)    at (0, 0)      {Gestionar requisitos\\de recuperación};
\node[caso] (uc3)    at (0, -1)     {Responder\\autoevaluación};
\node[caso] (uc4)    at (0, -2)     {Consultar propio\\perfil e historial};

%Casos de uso del administrador
\node[caso] (uc5)   at (0, -3.5)  {Gestionar usuarios\\y profesores};
\node[caso] (uc6)   at (0, -4.5)  {Administrar\\catálogos};
\node[caso] (uc7)   at (0, -5.5)  {Diseñar formularios\\de autoevaluación};
\node[caso] (uc8)   at (0, -6.5)  {Consultar cartas y\\requisitos de todos};
\node[caso] (uc9)   at (0, -7.5)  {Generar reportes\\de cumplimiento};

%Casos de uso del público
\node[caso] (uc10)  at (0, -9)    {Explorar UEA por\\periodo y programa};
\node[caso] (uc11)  at (0, -10)   {Ver carta o requisito\\publicado};

%Sistema
\begin{scope}[on background layer]
\node[sistema, fit=(login1)(uc1)(uc2)(uc3)(uc4)(uc5)(uc6)(uc7)(uc8)(uc9)(uc10)(uc11)] (sys) {};
\end{scope}
\node at (0, 2.9) {\textbf{Sistema CyAD}};

%Conexiones
\draw[conn] (docente) -- (login1);
\draw[conn] (docente) -- (uc1);
\draw[conn] (docente) -- (uc2);
\draw[conn] (docente) -- (uc3);
\draw[conn] (docente) -- (uc4);

\draw[conn] (admin) -- (login1);
\draw[conn] (admin) -- (uc5);
\draw[conn] (admin) -- (uc6);
\draw[conn] (admin) -- (uc7);
\draw[conn] (admin) -- (uc8);
\draw[conn] (admin) -- (uc9);

\draw[conn] (publico) -- (uc10);
\draw[conn] (publico) -- (uc11);

\end{tikzpicture}
